\documentclass[12pt,letterpaper,oneside]{article}
\newcommand{\planCaso}{AM Taller Autocentro}
\newcommand{\planDescripcion}{Borrador académico basado en documentación local; datos comerciales por confirmar.}
\newcommand{\planFecha}{2026-09-20}
\newcommand{\planLugar}{Cajeme, Sonora; no acredita domicilio comercial}
\input{ITESCA/maestria-en-gestion-administrativa/plan-de-negocios-mga/formato-itesca-plan-negocios-generadas.tex}
\begin{document}
\portadaITESCA{Mercadotecnia y publicidad}{6546}
\section{Etapa 1. Diagnóstico del encargo}

Este documento es un \textbf{borrador hipotético pendiente de confirmación}. Corresponde a la actividad 6546, ``Mercadotecnia y publicidad'', de la asignatura Plan de Negocios de la Maestría en Gestión Administrativa del ITESCA. El encargo requiere presentar un diseño publicitario y definir estrategias de introducción al mercado, canales de redes sociales y promociones viables.

El caso de referencia es \textbf{AM Taller Autocentro}, propuesta local de taller, llantera y refaccionaria. La identidad disponible comprende un emblema circular dorado, tipografía oscura y el lema ``Pasión por los autos. Compromiso contigo.'' Estos elementos se conservan sin modificar ni atribuir titularidad. No existe información suficiente para afirmar apertura comercial, domicilio autorizado, ventas, clientes, precios, inventario, contratos con proveedores o demanda validada. En consecuencia, las acciones se presentan como propuestas sujetas a comprobación.

El mercado objetivo preliminar estaría compuesto por personas que requieren mantenimiento mecánico, revisión de llantas o consulta de refacciones. Esto constituye un \textbf{supuesto de trabajo}; antes de difundir la publicidad deben confirmarse el territorio de atención, las necesidades del público, sus criterios de compra, sus hábitos digitales y la capacidad operativa del proyecto.

\section{Etapa 2. Planeación y verificación}

La preparación del producto se organiza en cuatro estados:

\begin{enumerate}
    \item \textbf{Investigación:} confirmar mercado objetivo, cobertura, servicios autorizados, horarios, presupuesto, canales de contacto, capacidad de atención y referencias de catálogo disponibles para cotización.
    \item \textbf{Generación:} convertir el flyer textual en una pieza gráfica que respete el emblema dorado, la tipografía oscura y el lema existente.
    \item \textbf{Revisión:} comprobar ortografía, legibilidad, jerarquía visual, coherencia comercial y ausencia de afirmaciones no documentadas.
    \item \textbf{Validación:} contrastar el resultado con la consigna, el formato de entrega y la rúbrica individual cuando estos datos sean proporcionados.
\end{enumerate}

Se considerarán errores bloqueantes la publicación de teléfonos, domicilios, precios o descuentos no confirmados; presentar referencias de catálogo como existencias; prometer resultados mecánicos sin diagnóstico; y afirmar que AM Taller Autocentro ya opera comercialmente. También debe evitarse habilitar canales que no tengan una persona responsable ni un proceso definido de respuesta.

\section{Etapa 3. Propuesta de mercadotecnia y publicidad}

\subsection{Objetivo de comunicación}

El objetivo es presentar a AM Taller Autocentro como una propuesta de atención automotriz que integra mecánica, llantas y consulta de refacciones mediante un proceso organizado de revisión, cotización y programación de citas. La comunicación facilitaría el primer contacto sin prometer disponibilidad, compatibilidad, tiempos de entrega o resultados que todavía no hayan sido comprobados.

\subsection{Flyer textual propuesto}

\textbf{AM Taller Autocentro}

\textbf{Mecánica, llantas y refacciones bajo revisión y cotización}

¿Tu automóvil necesita atención? Solicita información para programar una revisión de sus necesidades de mantenimiento, condición de las llantas o requerimientos de refacciones.

\begin{itemize}
    \item Orientación inicial sobre servicios de mecánica sujetos a diagnóstico.
    \item Consulta de opciones de llantas conforme a la medida previamente confirmada.
    \item Búsqueda de referencias de refacciones sujeta a compatibilidad y cotización.
    \item Programación propuesta de citas para organizar la atención.
\end{itemize}

\textbf{Pasión por los autos. Compromiso contigo.}

Solicita información por \textbf{[canal de contacto pendiente]}. Atención propuesta en \textbf{[horario por confirmar]}. La disponibilidad, compatibilidad, cotización y alcance de cada servicio deberán confirmarse antes de cualquier operación.

Para la composición visual se propone colocar el emblema circular dorado en la parte superior y, debajo, el nombre del proyecto con tipografía oscura y legible. El mensaje principal ocuparía la zona central, acompañado de una imagen autorizada relacionada con el cuidado automotriz. En la parte inferior se ubicarían el lema, el llamado a solicitar información y los datos de contacto confirmados. La pieza no deberá incorporar marcas de terceros, fotografías sin permiso, precios, existencias, descuentos ni un domicilio no autorizado.

\subsection{Estrategia de introducción al mercado}

La introducción se plantea en tres momentos. En la \textbf{preparación} se confirmarían los servicios que pueden comunicarse, la capacidad de atención, el procedimiento de citas, los responsables de los canales y los criterios de cotización. Esta fase evitaría que la publicidad genere solicitudes que el proyecto no pueda procesar.

En la \textbf{difusión inicial} se publicaría el flyer junto con contenidos que expliquen el proceso para solicitar información. La comunicación se concentraría en las categorías generales de mecánica, llantas y refacciones, sin presentar las referencias por cotizar como productos disponibles. La frecuencia de publicación deberá ajustarse al tiempo, presupuesto y personal realmente asignados.

En el \textbf{seguimiento} se registrarían las consultas, categorías solicitadas, citas propuestas, cotizaciones elaboradas y motivos de no atención. Esta información permitiría mejorar los mensajes y analizar el interés preliminar. Las consultas no deberán interpretarse automáticamente como ventas o evidencia de demanda validada.

\subsection{Canales digitales propuestos}

\begin{itemize}
    \item \textbf{Facebook:} difusión del flyer, información general, contenidos preventivos y avisos autorizados.
    \item \textbf{Instagram:} adaptación visual del flyer, publicaciones breves y materiales educativos sobre cuidado automotriz.
    \item \textbf{WhatsApp Business:} recepción organizada de consultas y seguimiento de citas, una vez confirmado el número institucional.
    \item \textbf{Perfil empresarial en mapas:} alternativa condicionada a la existencia de una ubicación pública y autorizada.
\end{itemize}

Todos los canales deberán mantener el mismo nombre, emblema, lema y tono informativo. Las respuestas deben indicar cuándo una solicitud requiere revisión técnica y evitar diagnósticos personalizados basados únicamente en mensajes o fotografías.

\subsection{Promociones y acciones complementarias}

Las promociones deberán aprobarse después de evaluar costos, capacidad y condiciones de aplicación. Se proponen las siguientes alternativas:

\begin{itemize}
    \item orientación inicial para identificar el tipo de revisión que podría requerirse;
    \item recordatorios voluntarios de citas o mantenimiento;
    \item paquetes de servicios relacionados, definidos únicamente después de calcular costos y alcance;
    \item publicaciones educativas estacionales sobre llantas y mantenimiento preventivo;
    \item colaboración con eventos locales o creadores de contenido solo cuando exista afinidad, autorización y un beneficio evaluable.
\end{itemize}

Los costos por estimar incluyen diseño, impresión, pauta digital, producción de contenido, atención de mensajes y seguimiento. Los posibles ingresos deberán clasificarse por servicios mecánicos, llantas y refacciones, sin asignar cifras hasta contar con costos, cotizaciones de proveedores y criterios de margen.

\subsection{Control y datos pendientes}

La evaluación propuesta consideraría alcance de publicaciones, consultas recibidas, solicitudes de cita, cotizaciones preparadas, tiempo de respuesta y motivos de no atención. Estos indicadores servirían para ajustar la estrategia, pero no equivaldrían por sí mismos a ventas ni a éxito comercial.

Para convertir este borrador en un trabajo real se requieren el formato y los límites de entrega, la rúbrica individual, el territorio, el mercado objetivo confirmado, el presupuesto, los responsables, los datos de contacto, los horarios, la capacidad operativa, el catálogo autorizado, las políticas de cotización y la evidencia de viabilidad de cada promoción.
\clearpage
\printbibliography[title={Fuentes de consulta}]
\end{document}
